\documentclass[11pt,a4paper]{article}
\usepackage[margin=2.2cm]{geometry}
\usepackage{booktabs,longtable,array}
\usepackage{siunitx}
\usepackage{hyperref}
\usepackage{url}
\newcommand{\nruns}{1564}
\newcommand{\nmodels}{9}
\newcommand{\nbenchmarks}{58}
\newcommand{\ncorrect}{1396}
\newcommand{\nbuildfail}{98}
\newcommand{\nselfcheckfail}{41}
\newcommand{\nruntimeout}{7}
\newcommand{\nagenterror}{22}
\newcommand{\nagenttimeout}{317}
\newcommand{\nsuspicious}{426}
\newcommand{\nslower}{152}
\newcommand{\medianspeedup}{213}
\newcommand{\suspiciouscap}{1000}


\title{LLM coding agents optimising serial HPC benchmarks on GB200:\\
       first results with OpenCode and RiVault-hosted models}
\author{Emil Vatai (harness and runs generated with Claude Code)}
\date{September 2026}

\begin{document}
\maketitle

\begin{abstract}
We let \nmodels{} open-weight LLMs, driven by the OpenCode coding agent, optimise
\nbenchmarks{} HeCBench benchmarks whose OpenMP pragmas had been stripped, on a
node with an NVIDIA GB200 GPU and a 144-core Grace CPU. Each (model, benchmark)
pair was attempted three times with a two-hour budget, for \nruns{} runs in total.
\ncorrect{} runs produced code that compiled, ran and passed the benchmark's
built-in self-check. The median verified speedup over the serial baseline is
\medianspeedup$\times$. A substantial fraction of reported speedups
(\nsuspicious{} runs above \suspiciouscap$\times$) are implausible and indicate
that the benchmarks' self-reported kernel timers can be defeated; these are
reported separately and excluded from the geometric means.
\end{abstract}

\section{Setup}

\paragraph{Benchmarks.}
HeCBench (branch \texttt{ktwork}, \texttt{src2/*-omp}) was reduced to its OpenMP
variants, every \texttt{\#pragma omp} line was deleted (leaving serial code that
still calls \texttt{omp\_*} runtime functions), and headers referenced from
sibling variants were copied in. Of 322 benchmarks, 284 compiled with
\texttt{nvc++} (NVIDIA HPC SDK 26.5), and \nbenchmarks{} both ran within
\SI{300}{s} and passed their built-in self-check (\texttt{VERIFY=yes}). Only
these \nbenchmarks{} were used, so correctness can be judged by the benchmark
itself.

\paragraph{Agents and models.}
OpenCode~1.18 with \nmodels{} models served by the RiVault LiteLLM gateway
(Table~\ref{tab:models}). Seven further \texttt{RiVault/*} aliases (Agentic,
Instruction and Reasoning in several sizes) were also run but are excluded from
this report because the underlying models are undisclosed and likely duplicate
the named ones. Claude Code runs were planned but not executed: the
available claude.ai subscription rate-limits at the scale required.

\paragraph{Protocol.}
Every run is one Slurm job (partition \texttt{gpu}, one GB200, \SI{2.5}{h}
limit) executing \texttt{harness/run\_agent.sh}:
\begin{enumerate}
  \item copy the benchmark twice; build and run the untouched copy to obtain the
        \emph{baseline} time;
  \item start the agent in the second copy with the fixed prompt of
        \texttt{harness/agent\_prompt.md} (optimise for this machine; any
        approach the compiler supports; do not touch verification or problem
        sizes; \SI{2}{h} limit); permissions are auto-approved; the agent is
        killed after \SI{2}{h};
  \item rebuild the agent's copy from clean, run it, require \texttt{PASS} in the
        output;
  \item speedup $=$ baseline time $/$ optimised time, where the time is the sum
        of the kernel/execution times the benchmark prints (wall time of the run
        if it prints none);
  \item record token usage and wall time; store an \texttt{fn-eval} feedback
        record with figure of merit \emph{speedup}.
\end{enumerate}
Each run had its own copy of the benchmark and its own OpenCode data directory;
runs never saw the repository, its history or the original OpenMP code.

\section{Results}

\subsection{Per model}

\begin{table}[h]
\centering\footnotesize\setlength{\tabcolsep}{4pt}
\begin{tabular}{l r r r r r r r}
\toprule
model & runs & correct & $>$\suspiciouscap$\times$ & median & geo-mean$^\dagger$ & wall & out.\ tok. \\
      &      &         &                          & speedup & speedup & (min) & (M) \\
\midrule
Seed-OSS-36B-Instruct & 174 & 92 & 2 & 7.45 & 5.53 & \textbf{120} & 3.7 \\
Qwen3.6-35B-A3B & 174 & 163 & 28 & 95.5 & 23 & 32.7 & 13.8 \\
Qwen3.8-27B & 174 & 143 & 68 & 866 & 28 & 120 & \textbf{19.3} \\
DeepSeek-V4-Flash & 173 & 163 & 33 & 189 & 43.4 & 8.07 & 7.4 \\
DeepSeek-V4-Flash-reasoning & 173 & 165 & 73 & 747 & 39.8 & 20.5 & 16.5 \\
Llama-4-Scout-17B-16E-Instruct & 174 & 171 & 0 & 1 & 1.07 & 0.216 & 0.1 \\
Kimi-K3 & 174 & 160 & 76 & 708 & \textbf{60.9} & 72.7 & 9.5 \\
GLM-5.2 & 174 & 158 & \textbf{79} & \textbf{997} & 27.8 & 111 & 16.8 \\
GLM-5.3-Flash & 174 & 163 & 67 & 478 & 36.8 & 52.2 & 18.0 \\
\bottomrule
\end{tabular}

\caption{Per-model summary over \nbenchmarks{} benchmarks $\times$ 3 repetitions
(provider prefixes dropped). ``correct'' = compiled, ran and passed the self-check;
``wall'' is the median agent wall time. Bold marks the column maximum in the
$>$\suspiciouscap$\times$ column and those to its right. $^\dagger$Geometric mean
over correct runs with speedup $\le$ \suspiciouscap$\times$; the median is over
all correct runs.}
\label{tab:models}
\end{table}

\begin{table}[h]
\centering\footnotesize\setlength{\tabcolsep}{4pt}
\begin{tabular}{l r l r r r r}
\toprule
model & peak & benchmark & \multicolumn{2}{c}{cost of the peak run} & \multicolumn{2}{c}{cost of the sweep} \\
      & speedup &         & tokens (k) & wall (min) & tokens (M) & agent-hours \\
\midrule
Seed-OSS-36B-Instruct & 894 & swish & 51.3 & 56.4 & 51 & 300 \\
Qwen3.6-35B-A3B & \textbf{997} & accuracy & 158 & 13.1 & 42.5 & 107 \\
Qwen3.8-27B & 945 & background-subtract & 281 & 120 & 46.5 & 287 \\
DeepSeek-V4-Flash & 990 & present & 150 & 13.3 & 17 & 35.3 \\
DeepSeek-V4-Flash-reasoning & 983 & floydwarshall & 201 & 20.7 & 41.8 & 72.6 \\
Llama-4-Scout-17B-16E-Instruct & 785 & mrc & 6.93 & 0.176 & 1.23 & 0.607 \\
Kimi-K3 & 853 & swish & 624 & 27.8 & 724 & 199 \\
GLM-5.2 & 975 & floydwarshall & 4.72e+03 & 120 & 1.28e+03 & 284 \\
GLM-5.3-Flash & 883 & swish & 156 & 32.9 & 54.1 & 159 \\
\bottomrule
\end{tabular}

\caption{Peak verified speedup per model (largest speedup $\le$
\suspiciouscap$\times$ among correct runs) and what it cost: tokens
(input$+$output) and agent wall time of that run, and the totals over the whole
sweep of \nbenchmarks{} benchmarks $\times$ 3 repetitions. The RiVault gateway
reports no monetary price, so cost is given in tokens and agent-hours.}
\label{tab:peaks}
\end{table}

\begin{table}[h]
\centering\footnotesize\setlength{\tabcolsep}{4pt}
\begin{tabular}{l r r r r r r}
\toprule
model & correct & build failed & self-check failed & run timeout & agent error & hit 2\,h \\
\midrule
Seed-OSS-36B-Instruct & 92 & 63 & 16 & 0 & 3 & 124 \\
Qwen3.6-35B-A3B & 167 & 4 & 1 & 0 & 2 & 0 \\
Qwen3.8-27B & 145 & 13 & 13 & 1 & 2 & 90 \\
DeepSeek-V4-Flash & 164 & 3 & 2 & 0 & 4 & 1 \\
DeepSeek-V4-Flash-reasoning & 166 & 0 & 0 & 0 & 7 & 2 \\
Llama-4-Scout-17B-16E-Instruct & 171 & 0 & 3 & 0 & 0 & 0 \\
Kimi-K3 & 163 & 2 & 3 & 2 & 4 & 23 \\
GLM-5.2 & 163 & 6 & 2 & 3 & 0 & 68 \\
GLM-5.3-Flash & 165 & 7 & 1 & 1 & 0 & 9 \\
\bottomrule
\end{tabular}

\caption{Run outcomes per model. ``agent error'' = the agent produced no output
tokens (API/server error). ``hit 2\,h'' counts runs in which the agent was
killed at the time limit (its last state is still verified and may be correct).}
\label{tab:outcomes}
\end{table}

Table~\ref{tab:peaks} lists each model's best verified result and its cost.
Three groups are visible in Table~\ref{tab:models}. Llama-4-Scout finishes
in seconds, changes essentially nothing, and therefore scores $\approx 1\times$
while being ``correct'' almost always. Seed-OSS-36B attempts changes but breaks
the build or the self-check in roughly half of the runs. The remaining models
(Kimi-K3, GLM-5.2, GLM-5.3-Flash, DeepSeek-V4-Flash and its reasoning variant,
Qwen3.8-27B, Qwen3.6-35B-A3B) keep 143--165 of 174 runs correct with median
speedups in the tens to hundreds.

\subsection{Per benchmark}

\begin{center}\small
\begin{longtable}{l r r r l r}
\toprule
benchmark & correct/runs & median & best$^\dagger$ & best model & $>$\suspiciouscap$\times$ \\
\midrule
\endfirsthead
\toprule
benchmark & correct/runs & median & best$^\dagger$ & best model & $>$\suspiciouscap$\times$ \\
\midrule
\endhead
\bottomrule
\endfoot
accuracy & 25/27 & 3.81e+03 & 997 & Qwen3.6-35B-A3B & 18 \\
ace & 21/27 & 150 & 590 & DeepSeek-V4-Flash & 0 \\
aligned-types & 24/27 & 189 & 566 & DeepSeek-V4-Flash-reasoning & 3 \\
atomicCost & 25/27 & 693 & 920 & Qwen3.8-27B & 6 \\
atomicPerf & 24/27 & 4.03e+04 & 344 & DeepSeek-V4-Flash-reasoning & 14 \\
background-subtract & 23/27 & 50.2 & 945 & Qwen3.8-27B & 1 \\
bezier-surface & 25/27 & 2.01 & 9.73 & Seed-OSS-36B-Instruct & 0 \\
bspline-vgh & 25/27 & 657 & 880 & Qwen3.6-35B-A3B & 11 \\
burger & 24/27 & 438 & 616 & DeepSeek-V4-Flash-reasoning & 1 \\
ccsd-trpdrv & 19/27 & 0.936 & 1.37 & Qwen3.8-27B & 0 \\
chacha20 & 26/27 & 688 & 866 & Qwen3.8-27B & 12 \\
colorwheel & 22/27 & 2.1e+03 & 161 & DeepSeek-V4-Flash & 13 \\
cross & 27/27 & 231 & 565 & DeepSeek-V4-Flash-reasoning & 0 \\
dp & 23/27 & 311 & 324 & GLM-5.3-Flash & 1 \\
floydwarshall & 23/27 & 564 & 983 & DeepSeek-V4-Flash-reasoning & 6 \\
fluidSim & 20/27 & 190 & 614 & Kimi-K3 & 0 \\
ga & 25/27 & 1.35e+03 & 805 & GLM-5.3-Flash & 15 \\
gaussian & 24/27 & 46.4 & 448 & Qwen3.8-27B & 2 \\
heat2d & 25/27 & 1.73 & 1.84 & DeepSeek-V4-Flash & 0 \\
inversek2j & 25/27 & 4.9e+03 & 685 & GLM-5.2 & 13 \\
ising & 24/27 & 1.68e+03 & 830 & Qwen3.6-35B-A3B & 18 \\
iso2dfd & 19/27 & 2.19 & 2.56 & DeepSeek-V4-Flash & 0 \\
jacobi & 26/27 & 146 & 413 & DeepSeek-V4-Flash & 1 \\
jenkins-hash & 26/27 & 1.36e+03 & 974 & DeepSeek-V4-Flash & 14 \\
kalman & 26/27 & 35.7 & 750 & DeepSeek-V4-Flash-reasoning & 0 \\
keccaktreehash & 19/27 & 2.21 & 63.2 & Kimi-K3 & 0 \\
layout & 25/27 & 5.48e+03 & 768 & Qwen3.6-35B-A3B & 16 \\
lif & 25/27 & 9.18e+03 & 309 & DeepSeek-V4-Flash & 16 \\
linearprobing & 18/27 & 1.05 & 1.24 & DeepSeek-V4-Flash-reasoning & 0 \\
lombscargle & 24/27 & 1.08e+04 & 113 & DeepSeek-V4-Flash & 18 \\
matrix-rotate & 14/27 & 19.6 & 918 & DeepSeek-V4-Flash-reasoning & 1 \\
maxpool3d & 26/27 & 328 & 522 & GLM-5.3-Flash & 2 \\
minisweep & 24/27 & 1.01 & 283 & Kimi-K3 & 0 \\
morphology & 23/27 & 0.00806 & 9.67 & Kimi-K3 & 0 \\
mr & 24/27 & 681 & 820 & DeepSeek-V4-Flash-reasoning & 10 \\
mrc & 25/27 & 185 & 785 & Llama-4-Scout-17B-16E-Instruct & 3 \\
nbody & 23/27 & 2.3e+03 & 584 & DeepSeek-V4-Flash-reasoning & 13 \\
nlll & 22/27 & 7.1 & 629 & DeepSeek-V4-Flash-reasoning & 2 \\
norm2 & 24/27 & 328 & 795 & GLM-5.2 & 3 \\
openmp & 25/27 & 27.7 & 443 & DeepSeek-V4-Flash & 2 \\
overlay & 21/27 & 2.52e+03 & 119 & Qwen3.6-35B-A3B & 16 \\
page-rank & 23/27 & 2.04e+03 & 676 & Qwen3.6-35B-A3B & 15 \\
perplexity & 24/27 & 927 & 955 & DeepSeek-V4-Flash-reasoning & 11 \\
present & 20/27 & 553 & 990 & DeepSeek-V4-Flash & 5 \\
projectile & 24/27 & 1.62e+03 & 334 & Seed-OSS-36B-Instruct & 14 \\
quantBnB & 27/27 & 6.83e+03 & 101 & Qwen3.6-35B-A3B & 21 \\
rainflow & 25/26 & 233 & 744 & Qwen3.8-27B & 5 \\
randomAccess & 25/27 & 25.1 & 681 & GLM-5.2 & 3 \\
rtm8 & 26/27 & 310 & 495 & DeepSeek-V4-Flash-reasoning & 6 \\
scel & 27/27 & 2.98e+03 & 771 & DeepSeek-V4-Flash & 18 \\
secp256k1 & 22/27 & 17.5 & 100 & Kimi-K3 & 0 \\
softmax & 27/27 & 2.1e+03 & 942 & DeepSeek-V4-Flash & 17 \\
stencil1d & 26/27 & 1.2e+03 & 894 & Qwen3.6-35B-A3B & 14 \\
swish & 25/27 & 894 & 894 & Seed-OSS-36B-Instruct & 12 \\
tensorT & 24/27 & 1.99e+03 & 257 & Qwen3.6-35B-A3B & 18 \\
triad & 24/26 & 0.887 & 1.61 & GLM-5.3-Flash & 0 \\
winograd & 25/27 & 1.45 & 152 & GLM-5.2 & 0 \\
zeropoint & 26/27 & 3.96e+03 & 712 & DeepSeek-V4-Flash & 16 \\
\caption{Per-benchmark results over all models and repetitions. $^\dagger$Best
verified speedup among runs $\le$ \suspiciouscap$\times$.}
\label{tab:benchmarks}
\end{longtable}

\end{center}

\section{Caveats}

\begin{itemize}
  \item \textbf{Timer gaming.} \nsuspicious{} correct runs report speedups above
        \suspiciouscap$\times$, up to $10^7\times$, and 25 report a kernel time
        of exactly zero. The self-checks passed, so the results are right, but
        the timed region evidently no longer contains the work (for example
        results computed once outside a repeated timing loop, or asynchronous
        GPU launches timed without synchronisation). The harness trusts the
        benchmark's own timer; an external, whole-program timing or an audit of
        the \texttt{work/} directories is required before these numbers are
        used. All aggregate speedups above exclude these runs.
  \item \textbf{Baseline.} The baseline is the stripped serial program on one
        Grace core, not the original OpenMP offload version. Speedups are
        therefore ``serial CPU $\to$ whatever the agent produced''.
  \item \textbf{Self-check strength.} Tolerances and coverage of the built-in
        checks vary; passing is necessary, not sufficient, for correctness.
  \item \textbf{Slower results.} \nslower{} correct runs were slower than the
        baseline; they count in the medians and geometric means.
  \item \textbf{Concurrency.} Up to 128 agents ran at once against one LiteLLM
        gateway; \nagenterror{} runs failed on the API side and \nagenttimeout{}
        hit the two-hour agent limit, both more likely under load.
\end{itemize}

\section{Reproduction}

\begin{verbatim}
./prepare_benchmarks.sh                 # HeCBench ktwork -> benchmarks/ (pragmas stripped)
./build_and_test_benchmarks.sh          # baseline report -> results/baseline_report.md
source .envrc && TOOLS=opencode harness/submit_all.sh   # 2784 Slurm array tasks
python3 harness/collect_results.py      # -> results/runs.json, results/summary.md
python3 report/make_tables.py && cd report && latexmk -pdf report.tex
\end{verbatim}
Per-run artefacts (agent transcript, modified sources, build and run logs,
\texttt{result.json}) are in
\texttt{../fn-agent-benchmarking-runs/opencode/<model>/<benchmark>/rep<N>/}.

\end{document}
